\chapter{Family--defect fermionic measure and Hubbard--Stratonovich origin gate}

Предыдущий audit оставил quantum-measure route как последнюю реализацию
middle action
\begin{equation}
 \tau_3(XX^T-|\Phi|^2I_3)^2.
\end{equation}
Здесь проверяется, возникает ли imaginary Hubbard--Stratonovich contour из
fermionic measure explicit KO6 geometry.

\section{Standalone finite KO6 bilinear}

Для linear part реальной структуры quadratic bilinear имеет matrix
\begin{equation}
 B_F=J_{F,\rm lin}D_F.
\end{equation}
В explicit 18-dimensional model машинно получено
\begin{equation}
 B_F^T=B_F.
\end{equation}
Его restriction на positive chirality имеет rank шесть, но antisymmetric
Grassmann part имеет rank ноль. Это согласуется с общей классификацией:
standalone real chiral fermion action в KO-dimension six обращается в нуль;
physical Pfaffian возникает после four-dimensional product с total
KO-dimension two. См. J.~W.~Barrett,
\emph{Fermion integrals for finite spectral triples},
\path{arXiv:2403.18428}.

\section{Completed Gaussian measure}

На radial slice
\begin{equation}
 X=\rho I_3,
 \qquad
 \Phi=r
\end{equation}
finite spectrum равен
\begin{equation}
 0^{\times6},
 \qquad
 \left(\pm\sqrt{\rho^2+r^2}\right)^{\times6}.
\end{equation}
После включения spacetime momentum \(p\) regulated modulus удовлетворяет
\begin{equation}
 \boxed{
 \log|\operatorname{Pf}\mathcal D_p|
 =3\log p^2+6\log(p^2+\rho^2+r^2).}
 \label{eq:family-final-pfaffian-sum}
\end{equation}
Все even moments имеют вид
\begin{equation}
 \Tr D_F^{2n}=12(\rho^2+r^2)^n.
 \label{eq:family-final-even-moments}
\end{equation}
Обе формулы проверены на 128 random points.

Точки
\begin{equation}
 (\rho,r)=(1,0),
 \qquad
 (\rho,r)=(2^{-1/2},2^{-1/2})
\end{equation}
имеют одинаковый Pfaffian при любом \(p\), но required moment-map square
равен на них соответственно
\begin{equation}
 1,qquad0.
\end{equation}
Следовательно, exact Gaussian measure не может породить required relative
curvature. Loop expansion даёт quartic ratio
\begin{equation}
 \rho^4:\rho^2r^2:r^4=1:2:1,
\end{equation}
а не \(1:-2:1\).

\begin{theorem}[Gaussian fermionic-measure no-go]
Spacetime completion делает fermionic integral ненулевым, но сохраняет
dependence только от \(\rho^2+r^2\). Поэтому ни exact Pfaffian, ни induced
spectral moments не воспроизводят
\((\rho^2-r^2)^2\).
\end{theorem}

\section{HS и torsion fork}

Imaginary Hubbard--Stratonovich identity корректно переписывает уже заданный
positive quartic, но не создаёт его. Current action не содержит выведенного
four-fermion channel на
\(\operatorname{Sym}_3(\mathbb R)=\mathbf1\oplus\mathbf5\), bare Gaussian
kernel или contour. Поэтому такой HS sector был бы новым parent input.

Spacetime Cartan/Holst torsion способен индуцировать four-fermion
interactions в иных theories; см. \path{arXiv:1004.1375} и
\path{arXiv:0911.5074}. Но project torsion является cohomological
\(\mathbb Z_2\)-line, а не dynamical Cartan field. Torsion-mediated rescue
меняет gravitational architecture и относится к следующей версии.

\begin{nogocriterion}[Pfaffian-to-HS no-go]
Нельзя объявлять imaginary HS identity следствием current Gaussian
Pfaffian. Для reopening требуется заранее выведенный four-fermion kernel,
modified differential calculus, Cartan torsion либо новый finite carrier.
\end{nogocriterion}

\section{Статус}

Family--defect route внутри Version IV закрыт. Tetrahedral projector,
three-cycle holonomy, gauge--family lock, KO6 embedding и moment-map
coefficient identity сохраняются как mathematical conditional structures,
но не объединяются current action в physical theory.

Машинный ledger сохранён в
\path{s2t_v4_family_defect_fermionic_measure_hs_gate_results.json};
генератор ---
\path{s2t_v4_family_defect_fermionic_measure_hs_gate.py}.